% Section 1.2 -- Assumptions, data collection & generation (max. 2 pages)
% Needs: booktabs, tabularx, amsmath. References: references.bib.
% Supporting material: app_data.tex (label app:data).

\section{Assumptions, data collection and generation}\label{sec:data}

The analysis gives a probability of failure $P_f$ and a reliability index
$\beta$ for each walmuur (the low wall between the werf and the canal) that
still stands on its original timber pile foundation. This needs, per wall:
whether the foundation has been renewed, the geometry of wall and foundation,
the retained height, the soil and water behind the wall, the load on the werf
and the condition of the timber piles. The municipal asset register lists the
walls, their length and age, but none of the structural data. We therefore
collected open data and project documents, derived the missing values with
scripts, and filled the remaining gaps with stated assumptions. Within the
time of the project, no data were requested from the municipality.

\subsection{Data collected}\label{sec:data-collected}

Table~\ref{tab:data-sources} lists the data, what each source provides and its
main weakness.

\begin{table}[htbp]
\centering
\small
\caption{Data used for the walmuren, with source and main limitation.}
\label{tab:data-sources}
\begin{tabularx}{\textwidth}{@{}>{\raggedright\arraybackslash}p{2.5cm}
  >{\raggedright\arraybackslash}p{4.3cm}
  >{\raggedright\arraybackslash}X >{\raggedright\arraybackslash}X@{}}
\toprule
Data need & Source & Used for & Limitation \\
\midrule
Walls, length, age & Asset register \cite{gu2025register} & 51 walmuren; screening & No height or foundation data \\
Werf level & AHN 0.5\,m \cite{ahn} & Retained height $H$ & Canal bed not measured \\
Soil & 71 DINO borings \cite{dino} & Soil class, sets $K$ and $\gamma$ & 20--144\,m away; often disagree \\
Groundwater & BRO wells \cite{bro_gmw} & GHG and GLG & Nearest well 143--418\,m \\
Traffic & BGT, NDW \cite{bgt,ndw} & Road distance, axle limit & Road is not on the werf \\
Foundation & Programme reports \cite{antea2019,gu2025up} & Renewed or timber & 8 walls unclear \\
Wall type & TNO dive notes \cite{tno2009} & $P(\text{type 2})$ per rak & Not known per wall \\
Section, piles & \cite{erfgoed156,erfgoed158,ibu2010} & Shapes, levels, piles & Schematic; Oudegracht \\
Timber, method & Hemel \cite{hemel2023,hemel2024} & Pile strength, decay, load & Amsterdam piles \\
Target $\beta$ & NEN 8700 \cite{nen8700} & 1.8 / 2.5 / 3.3 & Class not chosen \\
\bottomrule
\end{tabularx}
\end{table}

The foundation data decide which walls are modelled at all. Of the 51
walmuren, 32 are already renewed on steel. The register still gives 11 of
them the year of the old wall; counted from the renewal year, 18 walls instead
of 29 are past their design life. Of the 19 on timber, 2 have no werf:
the street runs to the water, which is a different structure. This leaves 17
walls (611\,m) for the $P_f$ calculation (Appendix~\ref{app:data}).

\subsection{Data generation}\label{sec:data-generated}

Eight scripts turn the raw data into one input row per wall. They run in a
fixed order, each reading the output of the one before
(Appendix~\ref{app:data}). The main generated values are:
\begin{itemize}\setlength{\itemsep}{0pt}\setlength{\parskip}{0pt}
  \item \textbf{Werf level and $H$:} the median of the AHN points on the werf
    side of the wall, leaving out points above +1.6\,m NAP on stairs or the
    street (+0.91 to +1.40\,m NAP). $H$ is the werf level minus the wall
    underside at $-0.08$\,m NAP \cite{erfgoed158}; median 1.21\,m.
  \item \textbf{Soil class:} the share of clay-like and sand-like soil between
    werf and wall underside in the nearest boring within 250\,m. At least
    70\,\% of one kind gives that class, otherwise ``mixed''. Of the 51 walls,
    23 are clay/peat, 16 sand and 12 mixed.
  \item \textbf{Foundation status:} a hand-compiled table of renewals per rak
    and bank, linked to the register through the city's \emph{Werfmuren} map.
  \item \textbf{Wall type:} divers found a timber beam (\emph{watersloof}, only
    in type 1) or a concrete floor (only in type 2) in some rakken
    \cite{tno2009}, which gives a chance of type 2 per rak.
  \item \textbf{Monte Carlo samples:} 400{,}000 draws per wall and wall type,
    with a fixed seed, for 2026 and 2041 (end of the 15-year reference period
    of NEN 8700).
\end{itemize}
Every raw file is listed with its source and download date, so all results
can be regenerated from the raw data.

\subsection{Assumptions}\label{sec:data-assumptions}

Where no data exist we use the assumptions in Table~\ref{tab:assumptions}.
Uncertain values get a distribution, not a single value. All 27 inputs are
listed in Appendix~\ref{app:data}.

\begin{table}[htbp]
\centering
\small
\caption{Main assumptions where data are missing.}
\label{tab:assumptions}
\begin{tabularx}{\textwidth}{@{}l X l@{}}
\toprule
Quantity & Assumption & Basis \\
\midrule
Canal bed & At the wall underside ($-0.08$\,m NAP), so $H$ is a lower bound & \cite{erfgoed158} \\
Wall type & $P(\text{type 2})$: 0.2 (rak 3--6), 0.5 (rak 21), 0.9 (rak 22), 0.8 (rest) & \cite{tno2009,erfgoed156} \\
Cross-section & Two measured drawings for all walls; floor level $\pm0.15$\,m & \cite{erfgoed156,ibu2010} \\
Pile condition & Soft shell $U(0,45)$\,mm, all walls; grows 0.10--0.25\,mm/yr & \cite{hemel2023} \\
Piles & Diameter $N(180\,\text{mm}, 0.10)$; spacing $U(0.8, 1.2)$\,m & \cite{ibu2010,hemel2023} \\
Soil & $K$ lognormal and $\gamma$ normal per class, independent & Group report, Tab.~C.5 \\
Water & Canal +0.6\,m NAP; groundwater $U(+0.5, +0.9)$\,m NAP & \cite{erfgoed158} \\
Surcharge & $q \sim N(7\,\text{kPa}, 0.20)$ at a distance $U(0,3)$\,m & \cite{hemel2024} \\
Proven strength & Keep only the draws in which the wall stands today & \cite{hemel2024} \\
\bottomrule
\end{tabularx}
\end{table}

\subsection{Effect on the results and validation}\label{sec:data-effect}

The assumptions do not affect the results equally. In order of effect on
$\beta$ (interpreted further in Section~4.2):
\begin{enumerate}\setlength{\itemsep}{0pt}\setlength{\parskip}{0pt}
  \item \textbf{Wall type.} For the same wall, type 1 gives a $\beta$ that is
    1.4--2.4 lower than type 2. With the chance per rak, $\beta$ is
    $-0.1$ to 1.5, so all 17 walls are below 1.8. If all were type 2, 8 of
    the 17 would pass 1.8. Whether a wall passes therefore depends mostly on its type, which the
    original drawings or a dive inspection can settle.
  \item \textbf{Pile condition.} For type 2 the piles fail far more often than
    the wall overturns (median $P_f$ 0.054 against 0.004). The same broad
    shell is used for every wall, so it shifts all $P_f$ but not the ranking.
  \item \textbf{Retained height and soil.} The soil moment grows with $H^3$:
    0.1\,m more height adds about 25\,\%. As $H$ is a lower bound, this
    assumption makes $P_f$ too low. Clay/peat gives 43\,\% more soil load than
    sand, and nearby borings agree for only 9 of the 51 walls.
  \item \textbf{Surcharge.} 7\,kPa is heavy for a pedestrian werf, so this
    assumption makes $P_f$ too high.
\end{enumerate}

We checked the data and assumptions where possible. The median AHN werf level
on the Oudegracht (+1.12\,m NAP) matches the +0.95 to +1.15\,m NAP measured in
\cite{erfgoed156,erfgoed158}. The renewed length in our foundation table
(about 3.36\,km) matches the city's 3{,}356\,m. The strongest check is that all
17 walls stand today. A first model that treated the wall as a masonry block
gave $P_f \approx 1$ and was rejected. With the measured sections, a type-2
wall stands in 95--100\,\% of the draws, a type-1 wall in only 33--98\,\%.
Either the type-1 model is conservative or many of these walls are type 2; the
proven-strength update accounts for this.

Settlement from pile rot, the damage actually seen on these walls
\cite{erfgoed156}, and failure of the timber screen are not modelled. The
$\beta$ values should therefore be read as an indicative ranking and order of
magnitude, not as exact values. The data that would improve them most are the
wall type per wall, the measured pile condition and canal soundings.
